Une bascule est formée d'une planche uniforme de \qty{4.0}{m} et de \qty{\v{M}}{kg} posée en son centre sur un pivot. Un enfant de \qty{\v{m1}}{kg} s'assoit à \qty{\v{d1}}{m} du pivot. Un second enfant, de \qty{\v{m2}}{kg}, s'assoit de l'autre côté.

\begin{enumerate}[label=\alph*)]
    \item À quelle distance du pivot le second enfant doit-il s'asseoir pour que la bascule soit en équilibre horizontal ?
    \item Quelle force le pivot exerce-t-il alors sur la planche ?
    \item Si le second enfant s'assoit plutôt à \qty{\v{d1}}{m} du pivot, quel est le module du moment de force résultant sur la bascule ?
\end{enumerate}
